\section{Bit-Ops and the Equivariance Lemma}
\label{sec:equivariance}

\subsection{Bit-ops}

\begin{definition}[Bit-op]
\label{def:bitop}
A \emph{bit-op} is an $\F_2$-linear endomorphism
$\bitop \in \End_{\F_2}(\cellring)$ that acts as a permutation
\[
  p_{\bitop} : \{0,\ldots,31\} \;\longrightarrow\; \{0,\ldots,31\}\,\cup\,\{\bot\}
\]
on coefficient positions, in the sense that
\[
  \bitop\!\Big(\sum_{j=0}^{31} c_j X^j\Big) \;=\; \sum_{j\,:\,p_{\bitop}(j)\ne\bot} c_j\, X^{p_{\bitop}(j)}.
\]
The two specific bit-ops considered in the implementation are, for $c \in \{0,\ldots,31\}$:
\begin{align*}
  \Rot_c &: \quad p_{\Rot_c}(j) = (j+c) \bmod 32, \\
  \ShR_c &: \quad p_{\ShR_c}(j) = j-c \text{ if } j \ge c, \text{ else } \bot.
\end{align*}
\end{definition}

\begin{remark}
\label{rem:lift}
A bit-op extends $\F_q$-linearly to $\liftring \to \liftring$ by acting
on coefficient positions in $\liftring$ exactly as on $\cellring$. We
denote the extension by the same symbol $\bitop$.
\end{remark}

\begin{definition}[Cell-wise action on columns]
For a bit-op $\bitop$ and a column $v \in \cellring^n$, define
$\bitop_*(v) \in \cellring^n$ by $(\bitop_*(v))_i := \bitop(v_i)$.
\end{definition}

\subsection{The equivariance lemma}

\begin{lemma}[Equivariance of MLE under bit-ops]
\label{lem:equivariance}
For any bit-op $\bitop \in \End_{\F_2}(\cellring)$, any column
$v \in \cellring^n$, and any $r \in \F_q^{\nu}$,
\begin{equation}
\label{eq:equivariance}
  \mle{\bitop_* v}(r) \;=\; \bitop\!\big(\mle{v}(r)\big) \qquad \text{in } \liftring.
\end{equation}
Equivalently, the diagram
\[
\begin{array}{ccc}
  \cellring^n & \xrightarrow{\;\bitop_*\;} & \cellring^n \\
  \downarrow{\scriptstyle \mle{(\,\cdot\,)}(r)} & & \downarrow{\scriptstyle \mle{(\,\cdot\,)}(r)} \\
  \liftring & \xrightarrow{\;\bitop\;} & \liftring
\end{array}
\]
commutes.
\end{lemma}

\begin{proof}
The MLE-evaluation map at a fixed point $r \in \F_q^{\nu}$ is
\[
  T_r : \cellring^n \;\longrightarrow\; \liftring,
  \qquad
  T_r(v) \;=\; \sum_{i \in \{0,1\}^{\nu}} \eq_i(r) \cdot v_i,
\]
with scalar coefficients $\eq_i(r) \in \F_q$. The map is
$\F_q$-linear in the $\liftring$-valued vector $(v_i)_i$, hence in
particular $\F_2$-linear in the cell coordinate $v_i \in \cellring$.

The bit-op $\bitop \in \End_{\F_2}(\cellring)$ acts trivially on the
row coordinate (it is the same map at every row $i$) and is
$\F_2$-linear in the cell coordinate. Since the row coordinate and
the cell coordinate of $\cellring^n = \cellring \otimes_{\F_2} \F_2^n$
are independent and the action of $\bitop$ commutes with scalar
multiplication by $\eq_i(r) \in \F_q$ (which lives in the row
coordinate after extension of scalars), we have
\begin{align*}
  T_r(\bitop_* v)
  &= \sum_{i} \eq_i(r) \cdot \bitop(v_i) \\
  &= \sum_{i} \bitop\!\big(\eq_i(r) \cdot v_i\big)
     \qquad \text{($\bitop$ is $\F_q$-linear after extension)} \\
  &= \bitop\!\Big(\sum_{i} \eq_i(r) \cdot v_i\Big)
     \qquad \text{($\bitop$ is additive)} \\
  &= \bitop\!\big(T_r(v)\big).
\end{align*}
\end{proof}

\subsection{Equivariance under projection}

\begin{corollary}[Projected equivariance]
\label{cor:projected-equivariance}
For any bit-op $\bitop$, any column $v \in \cellring^n$, any
$\alpha \in \F_q$, and any $r \in \F_q^{\nu}$,
\[
  \mle{\psia(\bitop_* v)}(r) \;=\; \psia\!\big(\bitop(\lmle{v}(r))\big).
\]
\end{corollary}

\begin{proof}
Both $\psia$ and $\bitop$ are $\F_q$-linear maps on $\liftring$, and
$\psia$ is also linear in the row coordinate (it acts on the cell
coordinate). By \Cref{lem:equivariance} applied to the column
$\psia(\bitop_* v) \in \F_q^n$, viewed cell-wise, and by linearity of
$\psia$,
\[
  \mle{\psia(\bitop_* v)}(r)
  \;=\; \psia\!\big(\mle{\bitop_* v}(r)\big)
  \;=\; \psia\!\big(\bitop(\mle{v}(r))\big)
  \;=\; \psia\!\big(\bitop(\lmle{v}(r))\big).
\]
\end{proof}

\Cref{cor:projected-equivariance} is the operative form: it tells the
verifier that the value of a virtual bit-op MLE at the final point
$r_0$ can be reconstructed from the lifted opening of the source
column by applying $\bitop$ to a polynomial in $\liftring$ and then
$\psia$, both of which are local field-arithmetic operations of cost
$O(32)$.

\begin{remark}[Geometric reading]
Under the identification $\cellring^n \cong \cellring \otimes_{\F_2}
\F_2^n$, $\bitop_*$ acts on the left tensor factor only and the MLE
evaluation $T_r$ is induced from a map on the right tensor factor
only. \Cref{lem:equivariance} is the assertion that the actions on
disjoint tensor factors commute, i.e. that a tensor product of
operators is well-defined.
\end{remark}

\subsection{Commutation with row-shifts}
\label{subsec:row-shift-commute}

Row-shifts act on the row coordinate. Fix $\delta \in \{0,\ldots,n-1\}$
and define $\rowsh{\delta} : \cellring^n \to \cellring^n$ by
\[
  (\rowsh{\delta} v)_i \;=\;
  \begin{cases} v_{i+\delta} & \text{if } i + \delta < n, \\ 0 & \text{otherwise,} \end{cases}
\]
where $i \in \{0,1\}^{\nu}$ is identified with its base-2 encoded
integer in $\{0,\ldots,n-1\}$ and arithmetic is over $\Z$.

\paragraph{Why row-shift is not a point shift.} It is tempting to seek
a transformation $\mathrm{shift}_\delta : \F_q^{\nu} \to \F_q^{\nu}$
satisfying $\mle{\rowsh{\delta} v}(r) = \mle{v}\big(\mathrm{shift}_\delta(r)\big)$
for all $v$. No such map exists in general: the integer-shift
permutation on $\{0,\ldots,n-1\}$ does not extend to an
$\F_q$-affine multilinear-shift compatible map on $\F_q^{\nu}$ (the
truncation at the top of the cube is the obstruction). The host PIOP
expresses the row-shift identity in \emph{kernel form} instead.

\begin{lemma}[Row-shift kernel identity]
\label{lem:row-shift-kernel}
There exists a bivariate polynomial
$\mathrm{next}_\delta(r, b)$, multilinear in $r \in \F_q^{\nu}$ for
each fixed $b \in \{0,1\}^{\nu}$, depending only on $\delta$ and not on
$v$, such that for every $v \in \cellring^n$ and every
$r \in \F_q^{\nu}$,
\begin{equation}
\label{eq:row-shift-mle}
  \mle{\rowsh{\delta} v}(r) \;=\; \sum_{b \in \{0,1\}^{\nu}} \mathrm{next}_\delta(r, b) \cdot v_b.
\end{equation}
A closed form on the boolean indexing $b \in \{0,1\}^{\nu}$ is
\[
  \mathrm{next}_\delta(r, b) \;=\;
  \begin{cases}
    \eq(r,\, b - \delta) & \text{if } b \ge \delta, \\
    0 & \text{if } b < \delta,
  \end{cases}
\]
extended multilinearly in $r$ to all of $\F_q^{\nu}$.
\end{lemma}

\begin{proof}
By definition of multilinear extension,
$\mle{\rowsh{\delta} v}(r) = \sum_{i \in \{0,1\}^{\nu}} \eq(r, i) \cdot (\rowsh{\delta} v)_i$.
Re-indexing $b := i + \delta$ and using
$(\rowsh{\delta} v)_i = 0$ for $i + \delta \ge n$,
\[
  \mle{\rowsh{\delta} v}(r) \;=\; \sum_{b \ge \delta} \eq(r, b - \delta) \cdot v_b
  \;=\; \sum_{b \in \{0,1\}^{\nu}} \mathrm{next}_\delta(r, b) \cdot v_b. \qedhere
\]
\end{proof}

\Cref{eq:row-shift-mle} expresses the row-shifted MLE as a fixed
$\F_q$-linear functional of the base column $v$, with weights
$\mathrm{next}_\delta(r, b)$ that depend on $\delta$ and $r$ but not on
$v$. The host PIOP's multi-point evaluation argument is built around
this kernel: it pairs the kernel as a selector MLE with the base
column MLE inside a single sumcheck (see \Cref{sec:mpeval}).

\begin{corollary}[Bit-op and row-shift commute]
\label{cor:bitop-shift-commute}
For any bit-op $\bitop$, any row-shift $\delta$, and any column
$v \in \cellring^n$,
\begin{equation}
\label{eq:bitop-shift-commute}
  \bitop_*\!\big(\rowsh{\delta} v\big) \;=\; \rowsh{\delta}\!\big(\bitop_*\, v\big) \quad\text{in } \cellring^n,
\end{equation}
and consequently
\begin{equation}
\label{eq:bitop-shift-mle}
  \mle{\bitop_*\!\big(\rowsh{\delta} v\big)}(r)
  \;=\;
  \sum_{b \in \{0,1\}^{\nu}} \mathrm{next}_\delta(r, b)\cdot \bitop(v_b).
\end{equation}
\end{corollary}

\begin{proof}
For every row $i$, both columns in \cref{eq:bitop-shift-commute} take
value $\bitop(v_{i+\delta})$ when $i + \delta < n$ and $0$ otherwise
(using $\bitop(0) = 0$ since $\bitop$ is $\F_2$-linear). For
\cref{eq:bitop-shift-mle}, apply \Cref{lem:equivariance} to
$\rowsh{\delta} v$:
\[
  \mle{\bitop_*(\rowsh{\delta} v)}(r) \;=\; \bitop\!\big(\mle{\rowsh{\delta} v}(r)\big)
  \;=\; \bitop\!\Big(\sum_b \mathrm{next}_\delta(r, b) \cdot v_b\Big),
\]
then use $\F_q$-linearity of $\bitop$ on $\liftring$ (\Cref{rem:lift})
to push $\bitop$ inside the sum.
\end{proof}

\Cref{cor:bitop-shift-commute} is the structural reason the
construction of \Cref{sec:construction} extends to row-shifted bit-op
virtual columns at no additional protocol cost: see
\Cref{sec:composition}.
